\documentclass{article}
\usepackage{iclr2026_conference,times}
\input{math_commands.tex}
\usepackage{hyperref}
\usepackage{url}
\usepackage{booktabs}
\usepackage{multirow}
\usepackage{graphicx}
\usepackage{microtype}

\title{When Does Sparse Conditional Capacity Help\\Scientific Image Adaptation Under Acquisition Shift?}
\author{Anonymous authors}

% Replace the ICLR 2026 style with the official 2027 release when available.
% Keep \iclrfinalcopy commented for anonymous submission.

\begin{document}
\maketitle
\begin{abstract}
Scientific images are collected across experiments, instruments, and clinical sites whose
acquisition signatures can change without changing the underlying scientific target. Sparse
mixture-of-experts (MoE) models offer conditional capacity that could absorb heterogeneous
adaptation demands, but the same mechanism could instead specialize to acquisition nuisance and
fail on unseen environments. We ask when sparse conditional capacity improves supervised
adaptation of a pretrained vision encoder under acquisition shift. After a predeclared dense-
substrate competence gate, we upcycle one feed-forward block of a dataset-qualified pretrained ViT
and compare it with an exactly initialized, total-parameter-matched dense block. A predeclared
$3\times2\times2\times3$ factorial varies placement,
routing unit, router geometry, and training pressure. Paired experiments separate conditionality
from parameter count, learned routing from random assignment, and route-level balance from
embedding-level invariance. Routing, counterfactual, and representation probes then test whether
successful models allocate capacity to biological content or acquisition environment. The study
is designed to produce operational guidance even when MoE does not win.
\end{abstract}

\section{Introduction}

Models for scientific images are routinely adapted across data collected by different experiments,
laboratories, scanners, or hospitals. These acquisition environments alter appearance while the
scientific prediction target remains nominally unchanged. A single dense feed-forward network must
represent every environment with the same parameters. Sparse mixture-of-experts (MoE) instead
supplies conditional capacity: each input activates only a subset of a larger parameter pool
\citep{fedus2022switch}. This
creates a plausible benefit---different adaptation demands need not interfere---and an equally
plausible failure---the router may partition data by acquisition nuisance and strand capacity when
a new environment appears.

Existing comparisons cannot answer which explanation is correct if an MoE is simply larger than
its dense baseline, if architecture and loss are tuned sequentially on the test domain, or if only
the winning configuration is analyzed. We therefore treat conditional capacity as the object of a
controlled kill test. We first require a pretrained transformer to demonstrate competent dense
full fine-tuning; a powerful model is not assumed to transfer merely because it is a foundation
model. We then change exactly one middle feed-forward block and hold the total parameter budget
fixed to the nearest realizable dense width. Only a seed-0 matched gain licenses replication or a
broader architectural study.

Our central question is: \emph{when adapting a pretrained vision model to scientific images
collected across acquisition environments, when does sparse conditional capacity improve
held-out-environment generalization over an equally large dense model?} We study RxRx1 cellular
microscopy, where 1,139 perturbations repeat across experimental batches and the official split
holds out complete experiments. The intended contributions are an auditable dense--sparse
comparison, a direct gradient-interference test, and mechanism analyses that distinguish reusable
conditional structure from acquisition nuisance.

\section{Scientific hypotheses}

\paragraph{H1: Conditional gain is not explained by parameter count.}
For a fixed total parameter budget, a learned top-1 MoE improves OOD performance only when routing
reduces interference among adaptation demands. The primary estimand is the paired difference
between MoE and a depth-matched dense-wide model, not MoE versus the original smaller encoder.

\paragraph{H2: Useful conditionality depends jointly on what can be routed and where.}
Image-level routing can exploit global acquisition context, whereas token-level routing can assign
local morphology independently of global appearance. Early layers retain low-level acquisition
statistics; later layers are closer to task semantics. We therefore expect routing unit and
placement to interact rather than admit an unconditional best choice.

\paragraph{H3: Router geometry and training pressure determine whether specialization follows
nuisance.}
Linear routing can use feature magnitude, which may encode staining, intensity, or contrast;
cosine routing removes this degree of freedom. Within-environment balancing constrains marginal
expert usage but does not remove environment information. Output-level adversarial training
directly discourages environment-decodable final embeddings.

\section{Controlled sparse adaptation}

\subsection{Pretrained substrate and upcycling}
We fully fine-tune a dataset-qualified pretrained ViT for supervised classification. Substrate
qualification is a Stage-0 gate: a candidate must reach a predeclared fraction of a canonical
validation-only sanity control before any MoE experiment is licensed. In each intervention,
exactly one transformer FFN is replaced. An MoE contains eight copies of the pretrained FFN and a
top-1 router; because every expert initially computes the same function, the MoE exactly matches
the qualified pretrained network at initialization. We place the converted block at early, middle,
or late depth. Exact parameter matching is recomputed separately for each qualified architecture.

\subsection{Total-parameter-matched dense control}
The primary dense comparator is one wider FFN at the same depth. Its integer hidden width is chosen
to minimize the total-parameter difference from the full MoE, including router parameters and
replicated expert biases. We do not add dormant padding parameters. We use an exact Net2Wider
Net2Wider construction \citep{chen2016net2net}: new hidden units copy source units, while each source unit's outgoing weight is split
unequally among its copies with coefficients that sum to one. The residual budget difference must
be below 0.1\%.

\subsection{Routing choices and pressures}
Image routing pools tokens before making one decision per image; token routing decides per patch.
Linear and cosine routers contain the same $E\times d$ weight matrix and one temperature scalar.
The \emph{canonical} condition uses global load balancing and router z-loss. The \emph{route}
condition computes load balance independently within each acquisition environment. The
\emph{output} condition retains canonical balancing and adds a gradient-reversal adversary
\citep{ganin2016domain} that
predicts training environment from the final embedding. Output invariance is applied identically
to MoE and dense controls.

\section{Experimental design}

\subsection{Datasets}
RxRx1 uses experimental batch as acquisition environment and siRNA perturbation as a 1,139-way
label. Camelyon17 uses hospital as environment and tumor presence as a binary label. We use the
official WILDS splits \citep{koh2021wilds}. Geometric
augmentation is shared across methods; photometric augmentation is excluded because it directly
manipulates the acquisition signal under study.

\subsection{Factorial and controls}
For each dataset, we run all $3\times2\times2\times3=36$ combinations of placement, routing unit,
router geometry, and pressure with three paired seeds. Dense-wide controls are run at all three
placements under canonical and output-invariance training. Route-pressure MoEs use the canonical
dense control because route-level balancing has no dense analogue.

\begin{table}[t]
\centering
\caption{Primary design. Every MoE cell is evaluated on both datasets with paired seeds.}
\label{tab:design}
\begin{tabular}{lll}
\toprule
Factor & Levels & Count\\
\midrule
Placement & early, middle, late & 3\\
Routing unit & image, token & 2\\
Router geometry & linear, cosine & 2\\
Training pressure & canonical, route, output & 3\\
\midrule
MoE configurations per dataset & & 36\\
\bottomrule
\end{tabular}
\end{table}

\subsection{Hyperparameters and stage gate}
We first tune one shared dense fine-tuning recipe using OOD validation. We then tune one router
temperature per dataset and geometry, one balancing coefficient per dataset, and one adversarial
coefficient per dataset. No cell receives its own search. Stage 1 completes the factorial without
evaluating OOD test. Stage 2 runs predeclared mechanism controls. Stage 3 evaluates a fixed
confirmatory set with fresh seeds and is the first stage allowed to access OOD test.

\subsection{Stage-0 substrate qualification}
Natural-image DINOv2 ViT-S/14 was the initial candidate on both datasets
\citep{oquab2023dinov2}. On RxRx1, a clean six-cell dense screen remained near floor, so a bounded
predeclared rescue tested removal of random-resized cropping, removal of layer-wise learning-rate
decay, their interaction, and official RxRx1 preprocessing. The best rescue used official
preprocessing and reached 5.50\% OOD-validation and 10.53\% seen-environment accuracy. A canonical
WILDS ResNet-50 sanity control reached 15.42\% best-so-far OOD-validation accuracy without
constructing OOD test. The rescue therefore achieved 35.7\% of the control, below the frozen 50\%
abandonment threshold. We exclude DINOv2 from the RxRx1 factorial and will qualify a microscopy-
pretrained ViT before running RxRx1 MoE cells. This is a substrate-validity decision, not an MoE
result. Camelyon17 qualification remains in progress. All reported Stage-0 selection uses OOD
validation, and OOD test remains untouched.

\subsection{Metrics and inference}
We report OOD accuracy, worst-environment accuracy, in-distribution accuracy, and the ID--OOD gap.
The primary contrast is the within-seed, within-placement difference between MoE and matched
dense-wide. A sum-to-zero factorial model estimates main effects and two-way interactions on
conditional gain. Seed-level paired bootstrap intervals quantify optimization uncertainty.

\section{Mechanistic analysis and result logic}

For every MoE we measure expert usage, assignment entropy, load imbalance, and dead experts. On a
held-out in-distribution audit set, we estimate acquisition-environment and biological-label
information in routes using mutual information and cross-validated decoders. We apply identical
linear probes to final embeddings to measure environment and label decodability.

We also replace learned assignments with seeded random assignments while freezing expert weights.
The accuracy change measures route reliance: a model can exhibit structured routing without
requiring that routing for its predictions. We relate conditional gain to route reliance,
environment decodability, label decodability, gradient conflict, and utilization across held-out
experiments. These are descriptive mechanism associations; causal claims are restricted to the
matched architectural contrast.

A positive matched gain with high route reliance supports useful conditional specialization. A
gain only against the smaller original encoder indicates capacity, not conditionality. High
acquisition decodability without matched gain indicates nuisance partitioning. If gradients
conflict but matched conditional gain remains absent across the held-out-experiment folds, then
sparse routing does not solve the interference mechanism proposed here.

\paragraph{Substrate qualification.}
Table~\ref{tab:native-competence} reports the predeclared native-stain Cell-DINO qualification.
Full fine-tuning recovers substantial task signal relative to a frozen linear readout, but its
performance drops by 14.80 points from seen experiments to OOD validation. The substrate is thus
competent enough to test conditional capacity, and the intended failure regime is batch transfer
rather than a floor representation. This result does not establish the proposed mechanism: the
same gap could arise from ordinary overfitting or acquisition variation that routing cannot fix.

\begin{table}[t]
  \centering
  \caption{Native-stain Cell-DINO qualification on RxRx1 (\%). OOD test is untouched.}
  \label{tab:native-competence}
  \begin{tabular}{lrrrr}
    \toprule
    Adaptation & Train & ID & OOD-val & Worst experiment \\
    \midrule
    Frozen backbone + linear head & 10.62 & 7.64 & 4.42 & 0.81 \\
    Full fine-tuning & 39.00 & 27.09 & 12.29 & 1.42 \\
    \bottomrule
  \end{tabular}
\end{table}

Table~\ref{tab:kill-contrast} reports the predeclared seed-0 kill contrast. Relative to the
exact-total-parameter-matched dense-wide model, MoE improves OOD-validation accuracy by 1.15 points,
ID accuracy by 1.98 points, and worst-experiment accuracy by 0.16 points. All four held-out
experiments move in the same direction. The gain is nevertheless far below the frozen 5-point
replication trigger, and the smaller original encoder remains 0.70 OOD-validation points better
than MoE. We therefore stop replication and architecture expansion. This single-seed gate does not
establish a zero population effect; it rejects the targeted 10--15-point effect strongly enough to
stop the predeclared campaign.

\begin{table}[t]
  \centering
  \caption{Seed-0 native-stain RxRx1 kill contrast (accuracy in \%). Dense-wide and MoE differ by
  0.001232\% in total parameters. OOD test is untouched.}
  \label{tab:kill-contrast}
  \begin{tabular}{lrrrrrr}
    \toprule
    Model & Params (M) & Active FFN (M) & Train & ID & OOD-val & Worst \\
    \midrule
    Original & 22.40 & 1.182 & 39.12 & 26.99 & 12.31 & 1.50 \\
    Dense-wide & 30.68 & 9.455 & 36.13 & 24.36 & 10.46 & 1.22 \\
    MoE & 30.68 & 1.185 & 38.72 & 26.35 & 11.61 & 1.38 \\
    \bottomrule
  \end{tabular}
\end{table}

The result does not support robustness from conditional capacity beyond ordinary capacity: MoE
recovers part of the degradation caused by dense widening but does not surpass the original
substrate.

The clean bounded routing diagnosis in Table~\ref{tab:route-diagnosis} further disfavors H2. The
learned router uses all eight experts, but its normalized usage entropy is essentially maximal and
route mutual information with experiment and perturbation is near zero. Randomizing its routes
reduces OOD-validation accuracy by only 0.42 points. A router frozen at initialization is 0.19
points better than the learned router and remains less than one point route-dependent. Learning the
router therefore does not identify a reusable batch-specialized partition that is necessary for
the model's predictions.

\begin{table}[t]
  \centering
  \scriptsize
  \caption{Clean seed-0 route diagnosis on OOD validation (accuracy and reliance in percentage
  points). Both controls use all eight experts. MI denotes the aligned token-route audit.}
  \label{tab:route-diagnosis}
  \begin{tabular}{lrrrrrr}
    \toprule
    Router & OOD-val & Random routes & Reliance & Entropy & MI(exp.) & MI(label) \\
    \midrule
    Learned & 11.65 & 11.23 & 0.42 & 0.9999 & 0.0051 & 0.0016 \\
    Frozen & 11.84 & 11.13 & 0.71 & 0.9975 & 0.0098 & 0.0041 \\
    \bottomrule
  \end{tabular}
\end{table}

This mechanism result is diagnostic rather than a population efficacy estimate. It is compatible
with sparse activation weakly regularizing a harmful dense widening, a fixed random partition
creating mildly useful subspaces, or single-seed variation. It does not establish the proposed
cross-batch gradient-interference mechanism, and it cannot reopen the failed replication gate. All
OOD-test values remain withheld.

\paragraph{Long-adaptation substrate-strength test.}
The bounded 90-epoch matrix in Table~\ref{tab:hypothesis90} separates a weak substrate from a weak
sparse intervention. Every adapted arm fits the training set perfectly and reaches 48--53\% ID
accuracy, while OOD-validation accuracy remains 18.5--20.5\%. Cell-DINO is therefore a competent
instrument and the remaining failure is transfer to unseen experiments, not inability to learn the
training task.

Canonical top-1 MoE improves OOD validation by 1.72 points and ID by 3.21 points relative to the
matched dense-wide model, but worst-experiment accuracy falls by 0.08 points. It is only 0.12 OOD
points above the smaller original model, which has a better worst-experiment score. The matched
models differ by 378 parameters (0.001232\%), within the predeclared 0.1\% tolerance. The primary
contrast therefore fails the five-point replication gate.

The matched OOD gain is $-0.42$, $+2.78$, $+2.55$, and $+1.72$ points at epochs 10, 30, 60, and
90, respectively. The conclusion is therefore already visible by epoch 60 and does not arise from
stopping before a growing sparse advantage; the final 30 epochs strengthen absolute baselines
while the conditional gain contracts.

\begin{table}[t]
  \centering
  \caption{Validated seed-0 90-epoch RxRx1 hypothesis matrix (accuracy in \%). Top-2 activates
  two experts and is not active-compute matched; partial and loss-based arms are diagnostic. OOD
  test is untouched.}
  \label{tab:hypothesis90}
  \begin{tabular}{lrrrr}
    \toprule
    Arm & Train & ID & OOD-val & Worst \\
    \midrule
    Original anchor & 100.00 & 52.09 & 20.09 & 1.75 \\
    Matched dense-wide & 100.00 & 48.23 & 18.50 & 1.70 \\
    Canonical token top-1 MoE & 100.00 & 51.45 & 20.22 & 1.62 \\
    Image top-1 MoE & 100.00 & 50.45 & 20.13 & 1.46 \\
    Within-experiment-balanced MoE & 100.00 & 51.96 & 20.46 & 1.38 \\
    Token top-2 MoE & 100.00 & 49.91 & 19.79 & 1.50 \\
    Frozen linear & 31.77 & 14.73 & 6.69 & 0.53 \\
    Last-four-block adaptation & 100.00 & 53.01 & 19.82 & 1.79 \\
    Output invariance & 100.00 & 43.53 & 16.58 & 1.38 \\
    Environment-balanced classification & 100.00 & 50.85 & 19.99 & 1.66 \\
    \bottomrule
  \end{tabular}
\end{table}

The alternative arms do not expose a hidden large sparse effect. Top-2 trails top-1, weakening
hard-routing starvation as the dominant explanation. Image routing is effectively tied with token
routing, so token granularity is not decisive. Within-experiment balancing gives the highest mean
OOD accuracy but the lowest worst-experiment score among the main adapted arms. Last-four-block
adaptation has the best ID and worst-experiment values without the best mean OOD value. These
patterns retain regularization, representation depth, and single-seed variation as alternatives;
they do not support a uniquely sparse or tail-robust mechanism. No replication, deeper mechanism
campaign, or OOD-test evaluation is opened by this result.

\paragraph{Exploratory post-gate confirmation.}
Under a subsequent explicitly multiplicity-exposed architecture search, the tail-safe seed-0
recipe was locked and compared with its exact-total-parameter-matched dense control at fresh seeds
1 and 2. At epoch 30, sparse minus dense changes OOD-validation/ID/worst-experiment accuracy by
$+1.76/+1.89/+0.04$ points for seed 1 and $+1.45/+1.50/+0.12$ points for seed 2. Thus the signs
align across both fresh seeds and all three decision axes, but the mean OOD effect remains far
below five points. These milestones are not a new tuned winner or a completed population claim:
all four locked arms continue unchanged to epoch 60, and the OOD test remains untouched. The
principal threat is selection across the preceding seed-0 factorial, router-auxiliary, and
temperature screens; contraction or sign reversal at epoch 60 is the predeclared falsifier.

\section{Limitations and broader relevance}

Two biomedical datasets do not establish a universal law of conditional adaptation. RxRx1 and
Camelyon17 were chosen because both expose real acquisition environments while differing sharply in
modality, number of environments, and label structure. Environment metadata is available during
training for two pressure conditions and for analysis, which may not hold in every deployment.
Adversarial invariance can also remove task-relevant variation when environment and label are
statistically coupled.

The dense-substrate gate may select different pretrained ViTs for microscopy and histopathology.
Consequently, architectural effects are identified within each dataset against its matched dense
control; cross-dataset agreement is evidence of replication across regimes, while differences
cannot be attributed to acquisition modality alone. We prefer this limitation to running a common
but floor-performing backbone that would make the MoE comparison uninterpretable.

The broader value is an auditable way to decide whether conditional capacity is appropriate for
heterogeneous scientific adaptation. The design separates added capacity, routing, and invariance;
the diagnostics reveal when experts follow nuisance; and negative results define when a simpler
dense adaptation is preferable.

\bibliography{references}
\bibliographystyle{iclr2026_conference}
\appendix
\section{Reproducibility checklist}

The released artifact will include exact dataset splits, configurations, seeds, parameter counts,
git commit, environment metadata, per-epoch logs, and per-environment sample counts. The aggregator
refuses to expose OOD-test metrics below Stage 3 and flags unmatched dense controls or budget
violations.

\section{Planned supplementary figures}

The appendix will contain expert-usage profiles, route and embedding decodability, counterfactual
rerouting, factor interactions, ID--OOD tradeoffs, and qualitative UMAP views. UMAP is
visualization only and is not used for hypothesis testing.

\end{document}
